\section{Proof of the main theorem}\label{sec:main-proof}

\begin{proposition}\label{prop:embedding}
Assume $\MC$. For every set $X$ there exist a well-orderable field $K$
of characteristic zero, a normal separable algebraic extension $L/K$,
and an injection $X\to L$.
\end{proposition}

\begin{proof}
For $X=\varnothing$, take $K=L=\Q$. Otherwise, by
Lemma~\ref{lem:partition}, fix a partition $(X_\alpha)_{\alpha<\lambda}$ of $X$ into nonempty finite blocks, and put $n_\alpha=|X_\alpha|$.
Let $\Q[t_x:x\in X]$ be the polynomial ring in distinct indeterminates $(t_x)_{x\in X}$ over $\Q$, and let $L$ be its field of fractions:
\[
 L=\Q(t_x:x\in X)
\]
Identify $\Q$ with its natural image in $L$. The map $x\mapsto t_x$ is
an injection from $X$ into $L$.

The family $(t_x)_{x\in X}$ need not be well-orderable, so we cannot use
these indeterminates directly to produce the required well-orderable base
field. Instead, for each $\alpha<\lambda$ we encode the unordered finite
block $\{t_u:u\in X_\alpha\}$, without enumerating it, by putting
\[
 p_\alpha(y)=\prod_{u\in X_\alpha}(y-t_u)
            =y^{n_\alpha}+\sum_{j<n_\alpha}c_{\alpha,j}y^j
            \in L[y].
\]
The coefficients are, up to sign, the elementary symmetric polynomials in
the variables $(t_u)_{u\in X_\alpha}$ and therefore do not depend on any
ordering of the block. In addition, the power $y^j$ canonically labels the
coefficient $c_{\alpha,j}$. Accordingly, put
\[
 D=\{(\alpha,j):\alpha<\lambda,\ j<n_\alpha\},
 \qquad \mathbf c=(c_{\alpha,j})_{(\alpha,j)\in D},
 \qquad K=\Q(c_{\alpha,j}:(\alpha,j)\in D).
\]
The lexicographic ordering $<_D$ is a well-ordering of $D$. Thus the
coefficients have canonical positions in a well-ordered family, without any
simultaneous choice of orderings of finite sets. Repetitions cause no
difficulty: a coefficient value that occurs more than once may be assigned
its least index in $D$.

Since $\Q$ is countable in $\ZFm$, fix a well-ordering $\prec_{\Q}$ of
$\Q$. For the extension tuple
$T=(L,\Q,\prec_{\Q},D,<_D,\mathbf c)$ we have $E_T=K$, so
Lemma~\ref{lem:orderExtension} gives a well-ordering of $K$.

It remains to recover the individual indeterminates algebraically. They need
not belong to $K$, but every $p_\alpha$ belongs to $K[y]$ and has precisely
the distinct roots $\{t_u:u\in X_\alpha\}$ in $L$. Hence, for every finite
$J\subseteq\lambda$, the field
\[
 E_J=K(t_x:x\in\bigcup_{\alpha\in J}X_\alpha)
\]
is the splitting field inside $L$ of
$\prod_{\alpha\in J}p_\alpha$. By Lemma~\ref{lem:splitting-normal},
it is finite and normal over $K$, and it is separable by
Lemma~\ref{lem:normal-separable}. These fields form a directed family whose
union is $L$. Consequently, every $a\in L$ lies in some $E_J$ and is
algebraic and separable over $K$. Moreover, the minimal polynomial of $a$
over $K$ splits in this $E_J$, because $E_J/K$ is normal. It therefore
splits in $L$, proving that $L/K$ is normal, separable and algebraic.
\end{proof}

\begin{proof}[Proof of Theorem~\ref{thm:main}]
By Theorem~\ref{thm:blass}, $\MC$ holds. Let $X$ be an
arbitrary set. We will show that $X$ is well-orderable.

By Proposition~\ref{prop:embedding}, there exist a well-orderable field $K$ of characteristic zero, a normal separable algebraic extension $L/K$, and an injection $X\to L$.
Thus, it suffices to show that $L$ is well-orderable.

For every well-orderable intermediate field $F$, the $F$-vector space
$L$ has a basis by hypothesis.
Thus, by Theorem~\ref{thm:algebraic}, $L$ is well-orderable and the proof is complete.
\end{proof}

With the same proof, we also have the following.

\begin{corollary}\label{cor:restricted}
In $\ZFm$, the following are equivalent:
\begin{enumerate}
\item $\AC$;
\item $\MC$+``every vector space over every well-orderable field of characteristic zero has a basis''.
\end{enumerate}
\end{corollary}
